%%%%%%%%%%%%%%%%%%%%%%%%%%%%%%%%%%%%%%%%%%%%%%%%%%%%%%%%%%%%%%%%%%%%%%%%%%
% Travtus CV -- Abhijeet Lokhande. AI Engineer (London Farringdon, hybrid).
% Travtus = "Everyday AI" platform for the housing industry. Platform-
% development AI role: shape how LLM-powered agents are designed,
% orchestrated, tested, monitored, improved in production. JD emphasis:
% agent architecture & multi-agent orchestration, tool use & function
% calling, API/service integration contracts, prompt tooling lifecycle
% (management, versioning, evaluation, testing, optimisation), model
% selection (quality/latency/scalability/cost), observability & guardrails,
% end-to-end/integration testing, POCs, solution design, performance/cost.
% Nice-to-have hook: prompt outputs used as code / code generation ->
% candidate's Spec-to-Code Agent. Honest: only cv.md facts. No front-end/UX
% claims. Font: Source Sans Pro. Teal + purple. Compile: pdflatex (2x).
%%%%%%%%%%%%%%%%%%%%%%%%%%%%%%%%%%%%%%%%%%%%%%%%%%%%%%%%%%%%%%%%%%%%%%%%%%

\documentclass[10pt,a4paper]{article}

\usepackage[T1]{fontenc}
\usepackage[utf8]{inputenc}
\usepackage[default]{sourcesanspro}   % professional, neutral sans (Source Sans Pro)
\usepackage[a4paper,top=0.7cm,bottom=0.5cm,left=1.4cm,right=1.4cm]{geometry}
\usepackage{titlesec}
\usepackage{enumitem}
\usepackage{xcolor}
\usepackage{hyperref}
\usepackage{microtype}

% ---- Colours (echo the HTML CV: teal primary, purple accent) ----
\definecolor{accent}{HTML}{117A8B}   % teal  (hsl 187,74%,32%)
\definecolor{company}{HTML}{6C2BB3}  % purple (hsl 270,70%,45%)
\definecolor{muted}{HTML}{555555}

% ---- Links: clickable label text, coloured teal (no raw URLs) ----
\hypersetup{colorlinks=true, urlcolor=accent, linkcolor=accent}

% ---- Remove paragraph indent, set sensible line spacing ----
\setlength{\parindent}{0pt}
\linespread{1.0}
\pagestyle{empty}

% ---- Section headings: uppercase, bold, teal, thin rule underneath ----
\titleformat{\section}
  {\normalfont\large\bfseries\color{accent}\scshape}
  {}{0em}{}[{\color{accent!35}\titlerule[1pt]}]
\titlespacing*{\section}{0pt}{3pt}{2pt}

% ---- Compact bullet list ----
\newlist{tight}{itemize}{1}
\setlist[tight]{leftmargin=1.3em,itemsep=1pt,topsep=1pt,parsep=0pt,label=\textcolor{accent}{\small$\bullet$}}

% ---- Entry: company / dates on line 1, role / location on line 2 ----
\newcommand{\entry}[4]{%
  {\color{company}\bfseries #1}\hfill {\color{muted}#2}\par
  \vspace{1pt}
  {\itshape #3}\hfill {\color{muted}\itshape #4}\par
  \vspace{3pt}%
}

% ---- Inline divider ----
\newcommand{\sep}{\,\textcolor{accent!60}{\textbar}\,}

\begin{document}

%======================= HEADER =======================
\begin{center}
  {\fontsize{24}{26}\selectfont\bfseries Abhijeet Lokhande}\\[3pt]
  {\color{accent}\large AI Engineer: LLM Agent Platforms, Orchestration \& Evaluation}\\[5pt]
  {\small
    London, United Kingdom \sep
    \href{mailto:abhijeet_lokhande@proton.me}{abhijeet\_lokhande@proton.me} \sep
    07480801382 \\[2pt]
    \href{https://www.linkedin.com/in/ablds/}{LinkedIn} \sep
    \href{https://github.com/abhijeetscode}{GitHub} \sep
    \href{https://huggingface.co/ab-ai}{Hugging Face}
  }
\end{center}
\vspace{2pt}

%======================= SUMMARY =======================
\section{Summary}
AI engineer who builds the platform capabilities that put LLM-powered agents into production: agent architecture, multi-agent orchestration, tool and function-calling patterns, and the prompt tooling lifecycle from management and versioning through evaluation, testing, and continuous optimisation. At a FinTech I have designed how agents interact with tools, APIs, and downstream services, defined the contracts those services expose, and kept it all reliable with observability, guardrails, and end-to-end testing. I make evidence-based model and framework choices on quality, latency, scalability, and cost, and I write production-grade Python. I work across experimentation, architecture, and delivery, and I open-source agents and models that other people use (67K+ HuggingFace downloads).

%======================= SKILLS =======================
\section{Skills}
\begin{tight}
  \item[\textcolor{accent}{\small$\bullet$}] \textbf{Agent architecture \& orchestration:} multi-agent orchestration (LangGraph, LangChain, LlamaIndex), tool use \& function calling, memory, workflow control, MCP, RAG
  \item[\textcolor{accent}{\small$\bullet$}] \textbf{Prompt \& LLM tooling lifecycle:} prompt engineering \& tuning, structured output design, prompt management \& versioning, evaluation \& testing tooling, continuous optimisation
  \item[\textcolor{accent}{\small$\bullet$}] \textbf{Evaluation, testing \& observability:} eval frameworks, LLM test strategies, end-to-end \& integration testing, tracing (OpenTelemetry), monitoring \& drift, guardrails \& failure handling, QoS \& outcome metrics
  \item[\textcolor{accent}{\small$\bullet$}] \textbf{Model \& performance:} model selection (quality, latency, scalability, cost), latency/throughput optimisation, cost trade-offs, POCs \& solution design
  \item[\textcolor{accent}{\small$\bullet$}] \textbf{Engineering:} production-grade Python (typed, tested, TDD), FastAPI, REST API design \& service contracts, SQL (PostgreSQL, MongoDB), Docker, CI/CD, AWS, GCP, Git
\end{tight}

%======================= EXPERIENCE =======================
\section{Experience}
\entry{Built AI (FinTech \& AI Company)}{April 2022 -- Present}{AI/ML Engineer (Early Stage Hire)}{London, UK}
\begin{tight}
  \item Designed agent architecture and orchestration: defined how agents interact with tools, APIs, and downstream services through a FastAPI Model Context Protocol (MCP) server exposing 180+ tools with standardized function-calling patterns, memory, and guardrails.
  \item Owned the prompt and LLM tooling lifecycle, from prompt design and structured outputs through evaluation, testing, and continuous optimisation of agent behaviour, iterating from real production usage.
  \item Kept agents production-ready with observability and testing: tracing (OpenTelemetry), monitoring, drift and incident alerts, groundedness checks, plus LLM test strategies and end-to-end integration testing across the platform.
  \item Made evidence-based model and framework choices on quality, latency, and cost, holding p95 around 175ms on a hybrid RAG service with a deterministic fallback, and delivered an 85\% performance gain by rearchitecting an inference engine into a graph-based pipeline.
  \item Built agents real users depend on (5+ leading investment firms) and rapid POCs that turned business and technical goals into durable platform capabilities, including an autonomous agent that cut analyst modeling time 70\%.
\end{tight}

\entry{King's College London}{March 2021 -- January 2022}{Research Associate}{London, UK}
\begin{tight}
  \item Improved drug-structure novelty 63\% by building deep learning models (Graph Convolution Networks, GANs), backed by a reproducible preprocessing pipeline over a million-row dataset.
\end{tight}

%======================= PROJECTS =======================
\section{Open-Source \& Projects}
\begin{tight}
  \item \textbf{Spec-to-Code Agent:} autonomous ReAct agent (LangGraph, Anthropic API) that uses structured prompt outputs as code, turning plain-language specs into installable, tested software packages; 9/9 pass rate across a polyglot benchmark (Python, Go, Rust, Java, JS). \href{https://github.com/abhijeetscode/Spec-to-Code-Agent}{View on GitHub}
  \item \textbf{Agentic Enterprise Assistant:} LLM agent over grounded data with RAG, role-based access at an MCP server (Keycloak JWT), a grounding-validation guardrail, auditable traces, and OpenTelemetry to Arize Phoenix. Passed 13/13 evaluation cases. \href{https://github.com/abhijeetscode/Agentic-Enterprise-Assistant}{View on GitHub}
  \item \textbf{Agentic Search API:} production retrieval service (FastAPI, Elasticsearch, Qwen3 embeddings, Claude tool-calling); an agent decomposes queries and routes to tools, with a deterministic fallback when the LLM fails. p95 around 175ms, 0.76 macro QoS score. \href{https://github.com/abhijeetscode/search-agent}{View on GitHub}
  \item \textbf{PII Detection LLMs:} fine-tuned GPT (LoRA) and BERT to detect personally identifiable information, deployed as reusable production artefacts; collectively 67K+ downloads on HuggingFace. \href{https://huggingface.co/ab-ai/pii_model}{View on Hugging Face}
\end{tight}

%======================= EDUCATION =======================
\section{Education}
{\color{company}\bfseries King's College London} \sep {\itshape MSc in Data Science}\hfill {\color{muted}London, UK}\par\vspace{1pt}
{\color{company}\bfseries C-DAC, India} \sep {\itshape Post Graduate Diploma in Advanced Computing (PG-DAC)}\hfill {\color{muted}Pune, India}\par\vspace{1pt}
{\color{company}\bfseries University of Pune} \sep {\itshape B.E. in Computer Science}\hfill {\color{muted}Pune, India}\par

\end{document}
